\subsection*{为何还需要另一种检查}

逐步终止服务能检验实际实现的回滚与恢复，却只执行了人为安排的若干交错。归纳论证可以覆盖任意有限步，但若允许步骤漏写了某个出口，或者保持理由中藏着错误，也可能在纸面上看似完整。可以再建立一个小模型，列出它的每项允许迁移，自动探索全部可达状态，并在每步判断不变量。

这种方法称为模型检查：把协议表示成状态与迁移，检查由这些迁移生成的执行是否满足给定性质。本章使用有界枚举，把身份、请求数量和输入种类限制为有限值，逐一探索交错。它适合找出短小反例；模型把数据库原子性等机制当作前提，检查在这些前提之上的协议安排，真实数据库与更大规模系统还须由各自的证据支持。

为了确保检查会发现目标错误，还要故意删除关键规则。若去掉身份唯一性以后检查依然全绿，可能是模型把两个请求排得过于顺序，或者压根没有统计同 \(S\) 的效果。错误变体的反例同时检验了模型能否表达那类失败，以及不变量有没有真正检查所关心的关系。

\subsection*{用两个身份和四次尝试表达竞争}

模型的两个逻辑身份为 S1、S2，除键 K1、K2 以外其他分量相同。四次物理尝试分别为 A0=(S1,原输入)、A1=(S1,原输入)、A2=(S2,原输入)、A3=(S1,另一输入)。原输入代表 \code{paper-001}、数量 1、Q7、CNY；另一输入把报价换为模型中的另一有效报价 Qm。两种模型输入都取 3000，差别作为独立标记保存，用来检查身份与输入比较。Qm 是这里的抽象报价，与第 4.3 节价格为 3500 的 Q8 分开；实际协议按原规范化版本逐字段核对报价与其他输入。

A0 与 A1 允许观察重复竞争、回滚后由另一尝试继续、已提交后响应丢失再恢复；A2 检查同参新身份能否独立形成记录；A3 检查已绑定时的异参冲突，也允许原尝试回滚后由另一输入首次绑定。物理尝试各自最多进入一次处理，失败后不会凭空复活；另一个尚未运行的尝试可以继续。于是恢复总次数有明确上限，这与任意多次自动重试有差别。

每次尝试有六个阶段：尚可开始、内部取得身份、准备业务且占有账户、已提交待交付、已交付结束、崩溃或响应丢失结束。模型保存当前余额、已提交终态集合、已提交效果集合、四个阶段，以及已经交付的业务结果。内部取得身份由尝试阶段表达，不写成外部已提交 \code{pending}。已交付结果进入状态，才能在后续账户变化后检查它与后来重放是否一致。

\begin{figure}[htbp]
\centering
\begin{tikzpicture}
\node[booknote,text width=11cm] at (5.6,4.4) {A0=(S1,Q7)，A1=(S1,Q7)，A2=(S2,Q7)，A3=(S1,Qm)};
\node[booknode,text width=2.5cm] (ready) at (1.4,3) {尚可开始};
\node[booknode,text width=2.5cm] (claim) at (5.4,3) {内部取得身份};
\node[booknode,text width=2.5cm] (prepare) at (9.4,3) {准备业务\\占有账户};
\node[bookresult,text width=2.5cm] (commit) at (9.4,.8) {已提交待交付};
\node[bookresult,text width=2.5cm] (deliver) at (5.4,.8) {已交付结束};
\node[booknode,text width=2.5cm] (drop) at (1.4,.8) {崩溃或丢失结束};
\draw[bookarrow] (ready) -- (claim);
\draw[bookarrow] (claim) -- (prepare);
\draw[bookarrow] (prepare) -- (commit);
\draw[bookarrow] (commit) -- (deliver);
\draw[bookalert] (claim.south) -- (5.4,2.1) -- (1.4,2.1) -- (drop.north);
\draw[bookalert] (prepare.south west) -- (7.4,1.6) -- (3.2,1.6) -- (drop.north east);
\draw[bookalert] (commit.south) -- (9.4,-.6) -- (1.4,-.6) -- (drop.south);
\node[booknote,fill=white] at (5.6,-.6) {提交后丢失响应：已提交事实保留};
\node[booknote,text width=10.5cm] at (5.6,-1.5) {已有终态时，尚可开始的尝试可直接重放或冲突结束\\每次物理尝试最多处理一次};
\end{tikzpicture}
\caption{四次物理尝试的六个阶段分别记录内部进度与交付结果}
\label{fig:ch06-model-phases}
\end{figure}


账户步骤在模型里一次只允许一个准备者占有，保持到提交或回滚。它抽象了本章单账户购买在条件判断与资金写入上的排队；拒绝也按这个简化次序处理。由于允许动作只减少余额，拒绝判断后不再足够的条件保持，简化没有借此覆盖充值。身份占用与账户占用分别检查，两个不同 \(S\) 可以先各自取得身份，再争账户。

模型没有展开订单与流水每个字段的写入，原子提交一次加入终态与关联效果，并记录当次余额。检查器另外实现前节的独立顺序对象：每条实现迁移发生时，根据原尝试 A0 等携带的身份与输入求预期，把实现投影为余额、终态映射和效果集合，逐项比较。首次提交必须产生规格要求的完整终态，内部步骤与回滚必须保持外部对象；重放和冲突还单独检查输出。独立规格不调用被检查迁移的余额决策代码，也不从新终态或流水反推身份、输入与金额。

这份抽象以尝试编号代表关联购买，提交一次发布终态与效果。因此它检查的是这个有限对象的迁移与输出；具体 SQL 中的 UUID、组合外键及结果字段，还要按第 6.1 节的方法从真实记录核对。

对于尚可开始的尝试，有原终态便比较输入、重放或冲突；无终态则检查同身份是否已被别的内部事务占有，可以取得时才进入下一阶段。内部事务可选择继续或崩溃回滚。完整准备后可原子提交，也可回滚。已提交以后可交付结果，也可丢失响应。等待没有改变状态，所以不反复增加同一个等待状态；只要竞争者可以结束，枚举就会探索它结束后的路径。

\subsection*{怎样知道已经枚举完当前边界}

从初始状态出发，每次取出一个尚未检查的状态，生成其中所有可执行的下一步。生成的完整状态若先前已经出现，就合并到同一节点，不再展开第二遍；若是新状态，加入待检查队列。这是状态去重，它比较余额、集合、所有尝试阶段和已经交付结果。两段同为 7000 的历史，可能一段有 S1 结果，另一段只有 S2，未来允许的恢复明显不同，因此必须保留为两个状态。

\begin{figure}[htbp]
\centering
\begin{tikzpicture}
\node[booknode,text width=4.5cm] (left) at (2.6,1.8) {状态甲\\余额 7000\\终态集合有 S1\\已交付 S1 的原结果};
\node[booknode,text width=4.5cm] (right) at (8.6,1.8) {状态乙\\余额 7000\\终态集合有 S2\\已交付 S2 的原结果};
\node[booknote] at (5.6,3.4) {余额相同，完整状态不同};
\node[bookresult,text width=4.5cm] (a) at (2.6,-.5) {后续 S1 同参请求\\重放原结果};
\node[bookresult,text width=4.5cm] (b) at (8.6,-.5) {后续 S1 同参请求\\可能首次取得身份};
\draw[bookarrow] (left) -- (a);
\draw[bookarrow] (right) -- (b);
\end{tikzpicture}
\caption{完整状态决定后继动作，余额相同的节点仍需分别保留}
\label{fig:ch06-fullstate}
\end{figure}


为了使反例容易复核，按距离初始状态的步数从短到长展开，称为广度优先搜索。每个首次发现的状态保留前一个状态与产生它的动作。某状态违反不变量时，沿这些前驱回到起点，便得到当前动作模型里步数最短的反例。这里的长度单位是一项抽象迁移；实际服务的耗时和 SQL 数量属于实现层面的测量指标。

\begin{figure}[htbp]
\centering
\begin{tikzpicture}
\node[booknode] (zero) at (.7,1.5) {初始};
\node[booknode] (a) at (3.3,2.4) {状态 A};
\node[booknode] (b) at (3.3,.6) {状态 B};
\node[booknode] (c) at (6.5,2.4) {状态 C};
\node[booknode] (d) at (6.5,.6) {状态 D};
\node[bookresult] (bad) at (10.1,2.4) {违反不变量};
\draw[bookarrow] (zero) -- (a);
\draw[bookarrow] (zero) -- (b);
\draw[bookarrow] (a) -- (c);
\draw[bookarrow] (b) -- (d);
\draw[bookarrow] (b) -- (c);
\draw[bookalert] (c) -- (bad);
\node[booknote] at (.7,-.4) {第 0 层};
\node[booknote] at (3.3,-.4) {第 1 层};
\node[booknote] at (6.5,-.4) {第 2 层};
\node[booknote] at (10.1,-.4) {第 3 层};
\node[booknote,text width=10.5cm] at (5.6,-1.3) {逐层展开，完整状态去重；首次发现时保存前驱与动作\\沿前驱回溯得到最短抽象迁移轨迹；队列清空表示当前有限图遍历结束};
\end{tikzpicture}
\caption{广度优先搜索按抽象动作数逐层展开并保留反例前驱}
\label{fig:ch06-bfs}
\end{figure}


本章分别以 10000、5000 与 2000 为初始余额运行正确规则。前两组各检查 1196 个不同状态、生成 2727 次迁移；2000 组检查 878 个状态、生成 2123 次迁移。每组最长的最短到达路径为 12 步；这是完整遍历当前有限图后的统计，并没有在第 12 步主动截断。三组分别覆盖两笔可成功、一笔成功后另一笔拒绝，以及从一开始就余额不足的数值分支。

\begin{table}[htbp]
\centering\small
\begin{tabularx}{\linewidth}{@{}rrrrX@{}}
\toprule
初始余额 & 不同可达状态 & 生成迁移 & 无后继状态 & 不变量反例\\
\midrule
10000 & 1196 & 2727 & 199 & 当前边界内未发现\\
5000 & 1196 & 2727 & 199 & 当前边界内未发现\\
2000 & 878 & 2123 & 115 & 当前边界内未发现\\
\bottomrule
\end{tabularx}
\caption{正确规则在两个身份、四次尝试边界内的完整枚举}
\end{table}

状态数量相同不表示业务结论相同。10000 与 5000 的阶段组合相同，成功与拒绝标签、余额和效果数量不同；在每组内部，完整状态都独立去重。没有后继也不等于所有逻辑操作都已完成：四次尝试可以全部在提交前崩溃，这也是允许的终止状态。有限图检查这些状态是否满足安全条件；最终完成所需的可用性、公平调度与持续恢复，仍按上一节的条件单独论证。

每条迁移同时接受结构不变量与独立顺序规格检查。前者检查守恒、身份数量、效果关联及原结果稳定；后者独立判断金额、资金足够时的成功、资金不足时的拒绝，以及同参恢复和异参冲突。检查器在生成迁移时同步核对这些关系，表中的状态仍按完整业务状态去重，检查代码本身没有增加新的业务阶段。

查询历史另行枚举一项购买和一项查询的全部 20 种交错，各自保持调用、提交或快照、返回的顺序。检查器按指定提交点与快照点逐项执行顺序对象，并核对点位于调用区间内及已完成后才开始的实际先后。跨超时的逻辑购买历史通过检查；成功返回后新权威查询仍报空的历史被拒绝。这是指定观察点的历史合法性检查，未把 20 条历史计入上表的数据库阶段状态图。

\subsection*{去掉身份唯一性，六步出现两次扣款}

第一个错误变体允许 A0 与 A1 同时内部取得 S1。账户排队仍在，原子提交也仍在，因此它并没有靠破坏余额条件制造异常。广度优先搜索在六步得到反例：A0 取得 S1；A0 准备成功；A1 也取得 S1；A0 提交到 7000；A1 准备成功；A1 提交到 4000。同 \(S\) 的两份终态已经成立。

\begin{figure}[htbp]
\centering
\begin{tikzpicture}
\node[booknode,text width=3cm] (a) at (1.6,2.6) {1. A0 取得 S1};
\node[booknode,text width=3cm] (b) at (5.6,2.6) {2. A0 准备成功};
\node[booknode,text width=3cm] (c) at (9.6,2.6) {3. A1 也取得 S1};
\node[bookresult,text width=3cm] (d) at (1.6,.4) {4. A0 提交\\余额 7000};
\node[booknode,text width=3cm] (e) at (5.6,.4) {5. A1 准备成功};
\node[bookresult,text width=3cm] (f) at (9.6,.4) {6. A1 提交\\余额 4000};
\draw[bookarrow] (a) -- (b);
\draw[bookarrow] (b) -- (c);
\draw[bookarrow] (c.south) -- (9.6,1.5) -- (1.6,1.5) -- (d.north);
\draw[bookarrow] (d) -- (e);
\draw[bookalert] (e) -- (f);
\node[booknote,text width=10.5cm] at (5.6,-1) {账户排队与余额条件仍成立；删除完整 S 排他后，同一身份形成两份成功};
\end{tikzpicture}
\caption{去掉身份唯一性后的六步反例}
\label{fig:ch06-no-unique}
\end{figure}


每笔扣款金额都正确，余额始终非负，效果与结果也各自原子保存。失败的是同 \(S\) 的数量上限。这个反例明确说明账户锁只能让两笔资金修改排队，无法单独决定它们属于同一购买。身份约束与账户保护观察不同对象，省去前者无法由后者补上。

独立检查在生成错误迁移时立即停止：已经展开 141 个状态、发现 294 个状态，检查 647 次迁移，得到六步反例。已发现与已展开的差别来自队列里尚未处理的节点。正确模型三组均在队列清空后统计，错误变体则只报告找到反例以前的搜索。

\subsection*{拆开提交，三步留下孤立效果}

第二个错误变体保留身份排他，却允许效果先提交、终态尚未保存。三步反例为 A0 取得 S1、准备扣款、提交 3000 的效果后退出。余额已经为 7000，效果集合有一条支出，终态集合为空，违反每条效果必须可由完成记录解释的条件。

这里甚至不必等到下一次重试真的重复扣款，状态已经丢失自动恢复需要的证据。后来的请求若把查空视为第一次购买，可能再扣一次；若选择永远不做，又可能让实际尚未执行的其他历史无法前进。共同提交把这项缺失状态排除在允许的发布步骤之外。

\begin{figure}[htbp]
\centering
\begin{tikzpicture}
\node[booknode,text width=3cm] (a) at (1.6,2) {1. 取得 S1};
\node[booknode,text width=3cm] (b) at (5.6,2) {2. 准备扣款};
\node[bookresult,text width=3cm] (c) at (9.6,2) {3. 先提交效果\\随后退出};
\draw[bookarrow] (a) -- (b);
\draw[bookalert] (b) -- (c);
\node[booknode,text width=10.5cm] at (5.6,-.1) {已提交余额 7000；效果集合有一笔 3000；终态集合为空\\扣款已经生效，原身份的完成记录尚未发布};
\end{tikzpicture}
\caption{拆开效果与终态提交，三步留下孤立扣款}
\label{fig:ch06-split-commit}
\end{figure}


第三个错误变体维持正确数据库写入，却在恢复成功结果时使用当前余额。七步反例为 S1 取得身份、准备、提交到 7000，S2 取得身份、准备、提交到 4000，随后 S1 同参恢复返回 4000。余额守恒、身份唯一与购买关联都可以通过，业务结果投影却与保存的原 7000 不同。恢复代码须从保存的原结果或 D1 读取当次 7000；当前账户的 4000 经独立查询展示，两种字段各自使用对应的数据来源。

\begin{figure}[htbp]
\centering
\begin{tikzpicture}
\node[bookresult,text width=3.1cm] (one) at (1.7,2.3) {1、2、3 S1 提交\\原结果当次 7000};
\node[bookresult,text width=3.1cm] (two) at (5.7,2.3) {4、5、6 S2 提交\\当前余额 4000};
\node[booknode,text width=3.1cm] (replay) at (9.7,2.3) {7. S1 同参恢复\\错误重算为 4000};
\draw[bookarrow] (one) -- (two);
\draw[bookalert] (two) -- (replay);
\node[booknode,text width=10.5cm] at (5.7,.1) {正确输出：从 S1 保存结果或 D1 取当次 7000\\当前账户 4000 另行查询\\历史结果与当前状态分别取值};
\end{tikzpicture}
\caption{写入全部正确，恢复时重算余额仍会破坏原结果投影}
\label{fig:ch06-rewrite-result}
\end{figure}


第四个错误变体把查找条件从完整 \(S\) 换为相同输入。S1 三步提交以后，S2 用同样商品与 Q7 进入，第四步就读回 S1 的结果。返回中购买所属身份与 A2 携带身份不同，违反不同身份分离的条件。即使这一次少扣钱而没有重复扣款，Alice 的第二项授权仍被吞掉；正确性检查必须同时观察过多与过少的业务操作。

\begin{figure}[htbp]
\centering
\begin{tikzpicture}
\node[bookresult,text width=4.5cm] (first) at (2.6,2.1) {1、2、3 S1 取得、准备、提交\\商品 paper-001，Q7，3000\\原购买 P1 属于 S1};
\node[booknode,text width=4.5cm] (second) at (8.6,2.1) {4. S2 使用相同输入进入\\错误按输入查找，返回 P1\\第二项确认被误合并};
\draw[bookalert] (first) -- (second);
\node[booknote] at (5.6,3.4) {错误按输入归组};
\node[booknote,text width=10.5cm] at (5.6,-.1) {正确恢复以完整 S 定位；输入比较用于核对该 S 的旧绑定\\新的 S2 应独立进入账户条件判断};
\end{tikzpicture}
\caption{按输入去重的四步反例吞掉不同身份的第二项购买}
\label{fig:ch06-input-dedupe}
\end{figure}


上述四种错误的最短轨迹分别为 6、3、7、4 步。还可以故意加入七种错误，检查独立规则是否确实承担了业务判断与输出检查的责任：

\begin{table}[htbp]
\centering\small
\begin{tabularx}{\linewidth}{@{}XrX@{}}
\toprule
错误迁移 & 最短步数 & 独立检查发现的违反\\
\midrule
余额足够却保存余额不足 & 3 & 应成功的业务决定错误\\
扣 4000 并写 4000 流水 & 3 & 守恒仍成立，报价金额错误\\
旧身份同参却报告冲突 & 4 & 输出不符合原输入绑定\\
旧身份再次计算并覆盖结果 & 4 & 恢复改变了已提交对象\\
新终态错写另一身份 & 3 & 与原尝试的 S 不符\\
新终态错写另一输入 & 3 & 与原尝试的输入不符\\
执行提交却遗漏终态 & 3 & 提交未产生完整规格迁移\\
\bottomrule
\end{tabularx}
\caption{错误变体对业务决定、投影与输出的独立检查}
\end{table}

这十一种错误全部找到反例。另一个查询变体在购买成功返回后开始权威查询，却返回未找到，也被历史检查拒绝。错误拒绝和错误金额仅凭结构不变量可以通过，独立规格能直接指出差别；错写身份、输入和漏终态的变体还检验了预期确实取自原调用，而没有跟随被检查输出一起改错。

\subsection*{把三种证据接回第一版协议}

归纳论证解释允许迁移怎样保持不变量；有界枚举展开规定身份与尝试数量内的全部交错，错误变体检验观察条件能否发现具体失败；第 6.1 节的中断方法则把状态与预期落实到真实数据库记录。模型以原子事务和已提交持久性为前提，实际终止检查应用层退出后的恢复，两者的故障范围都须与所用方法对应。增加请求数量、资金动作或恢复旧备份时，要相应扩展状态、允许迁移与检查依据。

对于当前购买，可以据这些规则准确描述一条允许历史：同 \(S\) 的多次尝试经过唯一仲裁，任一提交前退出只撤销本次内部组；一旦某次完整终态提交，后来同参请求恢复这份结果，新的 \(S\) 则按账户条件形成自己的成功或拒绝。客户端可能较晚才获知，记录必须在承诺窗口内继续可读。若下一步把扣款交给另一个系统，本地事务回滚将无法撤销对方已经完成的动作，效果与完成记录之间的共同提交关系就需要重新建立。
\FloatBarrier
